\section{Pares, financiamiento y complejidad organizacional}

Los capítulos anteriores trataron la cadena técnica y de producto; este capítulo añade las coordenadas externas de la empresa emergente. Gao Jiyang (高继扬) mencionó varias veces a sus pares durante la entrevista: aprender de Yushu (宇树) en el robot completo y la cadena de suministro, aprender de Physical Intelligence en modelos fundacionales y densidad de talento, y observar también lo que hacen empresas como Zhiyuan (智元) en gestión, propiedad intelectual y organización. Esta forma de expresarlo resulta interesante: no ve a sus pares solo como competidores, sino que trata a la industria como un sistema del que se puede aprender.

\subsection{Qué aprender de los pares}

Esta subsección desglosa «aprender de los pares» en fuentes de capacidad más concretas: las empresas de hardware ofrecen un modelo de cadena de suministro, las empresas de modelos ofrecen un modelo de inteligencia y los equipos directivos maduros ofrecen un modelo de organización. Vistos así, los competidores también son una fuente de conocimiento de la industria.

En el caso de Yushu, le interesa su profundidad en la cadena de suministro: capacidades de base como engranajes, motores, carcasas y simulación electromagnética. En el caso de Physical Intelligence, le interesan su posición de liderazgo en la dirección de la inteligencia, su densidad de talento y su densidad de capital. En el caso de Zhiyuan, le interesa cómo un equipo directivo maduro administra una empresa de inteligencia corporizada, incluidas la propiedad intelectual, los ajustes organizacionales y las decisiones de gestión pragmáticas.

\begin{knowledgebox}{Las tres dimensiones del aprendizaje de los pares}
\begin{tabularx}{\textwidth}{p{0.22\textwidth}p{0.32\textwidth}X}
\toprule
Referente & Qué se puede aprender & Relevancia para Xinghaitu (星海图) \\
\midrule
Yushu & Robot completo, cadena de suministro, desarrollo propio de componentes y profundidad de manufactura & Completa la base física de una empresa de robótica. \\
Physical Intelligence & Modelos fundacionales, densidad de talento e inversión en algoritmos de frontera & Referente para la capacidad aguas arriba del cerebro del robot. \\
Zhiyuan & Equipo directivo, propiedad intelectual, ajustes organizacionales y decisiones de operación & Aprender cómo una organización compleja entrega resultados de forma sostenida. \\
\bottomrule
\end{tabularx}
\end{knowledgebox}

\subsection{El aumento de la valuación no es la raíz de los problemas organizacionales}

En la entrevista se menciona que la valuación de Xinghaitu aumentó considerablemente en dos años y que el equipo pasó de poco más de diez personas a más de doscientas. Gao Jiyang sostiene que lo que realmente provoca los problemas organizacionales no es el aumento de la valuación en sí, sino que la organización se vuelve más compleja, que el scope de las tareas se amplía con rapidez, si los miembros originales y el equipo fundador pueden seguir el ritmo de crecimiento, y si es posible incorporar a tiempo a personas con más experiencia.

Las empresas de inteligencia corporizada enfrentan además un problema organizacional particular: la cadena de suministro del robot completo exige procesos, disciplina y calidad; el equipo de inteligencia exige densidad de talento, innovación y ensayo y error rápido. Estos dos domain tienen lenguajes de gestión distintos, ritmos distintos y formas distintas de tolerar errores. Reunirlos en una misma empresa es un desafío organizacional inherente al emprendimiento en inteligencia corporizada.

\begin{importantbox}{La doble naturaleza organizacional de las empresas de inteligencia corporizada}
De un lado está el hardware y la cadena de suministro, que requieren disciplina, procesos, calidad y control de costos; del otro, la inteligencia y los modelos, que requieren talento de alta densidad, exploración e iteración rápida. Poder gestionar ambas culturas al mismo tiempo es la cuestión central para que una empresa de robótica pueda crecer.
\end{importantbox}

\subsection{Resumen de la sección}

Xinghaitu no compite en el vacío. Por un lado aprende de sus pares; por el otro, atiende el financiamiento, la expansión organizacional y la gestión de dos domain. Para una empresa de inteligencia corporizada, que una ruta tecnológica llegue a cumplirse depende, en última instancia, de que la organización pueda soportar la complejidad.
